\section{Bài tập 4: Clamper Diode Circuit}

Cho mạch điện như hình bên dưới gồm nguồn một chiều $V$, diode $D_1$, điện trở $R_1 = 5.6\text{ k}\Omega$, diode ghim $D_2$ và điện trở tải $R_2 = 3.3\text{ k}\Omega$ mắc song song với $D_2$.

Giả sử sử dụng mô hình diode thực tế với điện áp phân cực thuận ngưỡng $V_D = 0.7\text{ V}$. Hãy thiết lập các phương trình và tính toán ba dòng điện $I_{R1}, I_{R2}, I_{D2}$ cùng điện áp trên tải $V_{R2}$ trong hai trường hợp điện áp nguồn:
\begin{itemize}[leftmargin=1.5em]
  \item Trường hợp 1: $V = 12\text{ V}$.
  \item Trường hợp 2: $V = 20\text{ V}$.
\end{itemize}
Sau đó, thực hiện mô phỏng PSpice ở chế độ Bias Point và rút ra nhận xét.

\given{ex4_given}{0.5\textwidth}{Mạch ghim điện áp sử dụng Diode (Hình 2.17)}

\subsection{Tính toán lý thuyết}
Khi cấp nguồn dương $V$, dòng điện chạy từ cực dương nguồn qua diode $D_1$. Do đó diode $D_1$ phân cực thuận và có độ sụt áp $V_{D1} = 0.7\text{ V}$.

\noindent Dòng điện đi tiếp vào nút song song giữa $D_2$ và $R_2$:
\begin{itemize}[leftmargin=1.5em]
  \item Diode $D_2$ có cực anode nối vào đầu trên và cực cathode nối vào đầu dưới của điện trở $R_2$.
  \item Do $V$ ($12\text{ V}$ hoặc $20\text{ V}$) đủ lớn, điện áp trên nhánh song song có xu hướng vượt quá $0.7\text{ V}$, do đó $D_2$ lập tức dẫn thuận và ghim điện áp trên nhánh tải ở mức:
  \begin{equation}
    V_{R2} = V_{D2} = 0.7\text{ V}
  \end{equation}
\end{itemize}

\noindent Dòng điện chạy qua điện trở tải $R_2$:
\begin{equation}
  I_{R2} = \frac{V_{R2}}{R_2} = \frac{0.7\text{ V}}{3300\,\Omega} \approx 0.2121\text{ mA}
\end{equation}

\noindent Áp dụng định luật KVL cho toàn vòng mạch chính:
\begin{equation}
  V - V_{D1} - V_{R2} - I_{R1} \cdot R_1 = 0 \iff V - 0.7 - 0.7 - I_{R1} \cdot R_1 = 0
\end{equation}
\begin{equation}
  I_{R1} = \frac{V - 1.4\text{ V}}{R_1} = \frac{V - 1.4}{5600}
\end{equation}

\noindent Áp dụng định luật KCL tại nút phân nhánh:
\begin{equation}
  I_{R1} = I_{D2} + I_{R2} \implies I_{D2} = I_{R1} - I_{R2}
\end{equation}

\subsubsection*{1. Khi $V = 12\text{ V}$:}
\begin{align*}
  I_{R1} &= \frac{12 - 1.4}{5600\,\Omega} = \frac{10.6}{5600} \approx 1.8929\text{ mA} \\
  V_{R2} &= 0.7000\text{ V} \\
  I_{R2} &= \frac{0.7}{3300\,\Omega} \approx 0.2121\text{ mA} \\
  I_{D2} &= I_{R1} - I_{R2} = 1.8929 - 0.2121 = 1.6808\text{ mA}
\end{align*}

\subsubsection*{2. Khi $V = 20\text{ V}$:}
\begin{align*}
  I_{R1} &= \frac{20 - 1.4}{5600\,\Omega} = \frac{18.6}{5600} \approx 3.3214\text{ mA} \\
  V_{R2} &= 0.7000\text{ V} \\
  I_{R2} &= \frac{0.7}{3300\,\Omega} \approx 0.2121\text{ mA} \\
  I_{D2} &= I_{R1} - I_{R2} = 3.3214 - 0.2121 = 3.1093\text{ mA}
\end{align*}

\begin{dapsobox}
  \begin{itemize}[leftmargin=1.5em]
    \item Khi $V = 12\text{ V}$: $I_{R1} = 1.893\text{ mA}$, $I_{R2} = 0.212\text{ mA}$, $I_{D2} = 1.681\text{ mA}$, $V_{R2} = 0.700\text{ V}$.
    \item Khi $V = 20\text{ V}$: $I_{R1} = 3.321\text{ mA}$, $I_{R2} = 0.212\text{ mA}$, $I_{D2} = 3.109\text{ mA}$, $V_{R2} = 0.700\text{ V}$.
  \end{itemize}
\end{dapsobox}

\Needspace{0.66\textheight}
\subsection{Mô phỏng PSpice}
\simresult
\simfigs{ex4_sim1}{Kết quả mô phỏng Bias Point của Bài tập 4 với $V = 12\text{ V}$}
\simfigs{ex4_sim2}{Kết quả mô phỏng Bias Point của Bài tập 4 với $V = 20\text{ V}$}

\subsection{So sánh và nhận xét}
\begin{table}[H]
  \centering
  \caption{So sánh kết quả tính toán lý thuyết và mô phỏng PSpice cho Bài tập 4}
  \small
  \setlength{\tabcolsep}{3.5pt}
  \begin{tabular}{lcccccccc}
    \toprule
    & \multicolumn{4}{c}{\textbf{Lý thuyết (Theory)}} & \multicolumn{4}{c}{\textbf{Mô phỏng (PSpice)}} \\
    \cmidrule(lr){2-5}\cmidrule(lr){6-9}
    \textbf{Nguồn} & $I_{R1}$ & $I_{R2}$ & $I_{D2}$ & $V_{R2}$ & $I_{R1}$ & $I_{R2}$ & $I_{D2}$ & $V_{R2}$ \\
    & {\scriptsize mA} & {\scriptsize mA} & {\scriptsize mA} & {\scriptsize V} & {\scriptsize mA} & {\scriptsize mA} & {\scriptsize mA} & {\scriptsize V} \\
    \midrule
    $V = 12\text{ V}$ & 1.893 & 0.212 & 1.681 & 0.700 & 1.903 & 0.203 & 1.701 & 0.669 \\
    $V = 20\text{ V}$ & 3.321 & 0.212 & 3.109 & 0.700 & 3.327 & 0.208 & 3.119 & 0.684 \\
    \bottomrule
  \end{tabular}
\end{table}

\begin{notebox}{Kết luận từ bảng}
\begin{itemize}[leftmargin=1.5em,itemsep=2pt,topsep=0pt]
  \item Kết quả mô phỏng khớp tốt với lý thuyết: $I_{R1}$ và $I_{D2}$ chênh dưới $1.5\%$; $I_{R2}$ và $V_{R2}$ thấp hơn lý thuyết khoảng $2\%$ đến $4\%$ do diode mô phỏng sụt áp thuận $0.67\text{ V}$ đến $0.68\text{ V}$ thay vì $0.7\text{ V}$.
  \item Khi nguồn tăng từ $12\text{ V}$ lên $20\text{ V}$, điện áp trên tải $V_{R2}$ gần như không đổi ($0.669\text{ V}$ lên $0.684\text{ V}$) vì bị $D_2$ ghim, nên dòng tải $I_{R2}$ cũng gần như không đổi ($\approx 0.2\text{ mA}$).
  \item Phần điện áp tăng thêm rơi trên $R_1$ và đi qua $D_2$: $I_{R1}$ tăng từ $1.903\text{ mA}$ lên $3.327\text{ mA}$, $I_{D2}$ tăng từ $1.701\text{ mA}$ lên $3.119\text{ mA}$.
\end{itemize}
\end{notebox}
